\section{Introduction}

The previous chapter defined PatchOps as a risk-aware decision system for failed CI workflows. This chapter
shows how that definition is realized architecturally. The most important high-level point is that PatchOps is
not currently a free-form multi-agent swarm. It is a structured, multi-stage pipeline in which only some stages
use LLM calls and the control flow itself remains explicit. This is a deliberate design decision, not a missing
feature.

The architecture follows the evidence discussed in the literature review. Repair systems are easier to defend,
measure, and ablate when the stages and transitions are visible. PatchOps therefore separates deterministic
preprocessing, LLM-backed reasoning, deterministic policy logic, and downstream validation into distinct steps.
This separation gives the thesis a clear systems story: repository-aware repair is not one monolithic prompt,
but a controlled workflow that transforms raw CI evidence into an operational decision.

\section{Architectural Philosophy}

PatchOps is best understood as a \emph{multi-stage decision pipeline} with an implemented validation harness
around selected candidate patches. The core diagnosis-to-decision engine is single-pass, and the stronger
multi-attempt verifier/rediagnosis loop remains the target architecture. The system already includes several
LLM-backed components, but they are embedded in a deterministic state machine rather than being allowed to
coordinate autonomously. In the terminology used throughout the project:

\begin{enumerate}
\item a \textbf{stage} is one step in the workflow,
\item a \textbf{deterministic stage} uses fixed rules or thresholds and does not call an LLM,
\item an \textbf{LLM-backed stage} performs bounded model-assisted reasoning, and
\item an \textbf{agentic loop} is a repeated act-observe-revise process in which the result of one action feeds
back into another reasoning step.
\end{enumerate}

This distinction is important because it clarifies what is already built and what is still proposed. The
implemented PatchOps engine is a single pass from log to decision, and the project includes a local Docker
validation harness for checking selected candidate fixes. The target architecture adds a tighter
generate-validate-rediagnose loop and a verifier stage, which are the pieces that make the workflow genuinely
agentic in the stronger sense.

\section{Current Implemented Workflow}

The current PatchOps workflow begins with a failed CI log and resolves to exactly one final action. The
implemented stage order is:

\begin{enumerate}
\item log extraction,
\item failure classification,
\item flaky or transient detection,
\item repository context retrieval,
\item root-cause analysis,
\item patch planning and intended-action selection,
\item semantic ambiguity checking,
\item patch generation when a code change is warranted,
\item safety gating,
\item plausibility gating for candidate automatic patches, and
\item final decision emission.
\end{enumerate}

This design directly matches the system objective established in Chapter 3. The pipeline does not attempt to
``think about everything at once.'' Instead, it narrows the decision space incrementally. Raw logs are first
compressed into evidence. Evidence is then classified. The classification and extracted signals are used to
decide whether the case is flaky, low risk, medium risk, high risk, or ambiguous. Only after these earlier
filters does the system spend effort on candidate patch generation.

In the refined Phase 15 architecture, the low-risk act-versus-hold branch can also be routed through a learned
decision gate. This gate does not replace the deterministic routing for flaky, high-risk, or unknown cases.
Instead, it replaces the earlier hand-set high-confidence threshold on the low-risk lane with a calibrated model
that predicts whether PatchOps should act or hold before generation.

\section{Stage 1: Log Extraction}

The first stage is a deterministic log extractor that performs the same kind of evidence compression suggested by
the CI diagnosis literature. Its responsibilities are to remove formatting noise, identify the failed step,
extract a critical error block, and detect known failure signatures. This stage is crucial because CI logs are
far too noisy to be treated as raw model context in every case.

Architecturally, this stage serves two purposes. First, it reduces token cost and prompt clutter. Second, it
creates structured evidence that can be reused by several later stages rather than forcing each one to reparse
the original log independently. This is an example of design discipline that matters in systems research:
expensive reasoning is deferred until the most relevant evidence has already been isolated.

\section{Stage 2: Failure Classification}

The classification stage maps the compressed failure evidence into the project's operational taxonomy. In the
current architecture, this stage is LLM-primary but can also incorporate rule hints and other supporting
signals. The output is not only a category label but also a confidence value and rationale. That confidence
value becomes a control signal for later policy decisions, which is why classification is not merely a reporting
step.

This stage exists because subsequent decisions depend on failure type. A runtime version mismatch, a flaky
network timeout, a missing secret, and a failing test should not be pushed through the same patching logic.
Classification therefore acts as the bridge between raw diagnosis and risk-aware action selection.

\section{Stage 3: Flaky and Transient Detection}

Once the system has a compressed view of the failure, it checks whether the run looks flaky or transient. This
stage is deterministic in the current implementation and is intentionally placed before repository reasoning and
patch generation. The rationale is straightforward: if the observed failure is likely due to a network timeout,
registry instability, connection reset, or similar transient issue, the safest response is rerun rather than
repair.

This ordering is one of the clearest expressions of the thesis contribution. Many repair systems ask ``can I
patch this?'' before asking ``should I patch this?'' PatchOps flips that priority by giving transient handling an
early exit branch in the architecture.

\section{Stage 4: Repository Context Retrieval}

If the failure is not classified as transient, PatchOps retrieves repository context. This stage is a key
architectural differentiator from log-only diagnosis assistants. The retriever collects category-relevant files
such as workflow YAML, manifests, configuration files, lockfiles, and evidence-linked source files. In the
expanded full-repository mode described later in the thesis, it can also surface candidate fix files that the
log never names directly.

The architectural role of this stage is to make the later reasoning repository-aware without letting repository
size overwhelm the system. PatchOps does not dump the full repository into the model context. Instead, it uses a
bounded retrieval step to select files that are likely to matter for the current failure. This is one of the
main ways the system balances practicality with context sensitivity.

\section{Stage 5: Root-Cause Analysis}

The root-cause analysis stage is LLM-backed. It takes the extracted evidence, the classification result, and the
retrieved repository files, and produces a summarized explanation of the likely defect and the affected files.
Architecturally, this stage converts local evidence into a hypothesis about what actually went wrong.

This stage is distinct from classification for a reason. Classification says what type of failure occurred.
Root-cause analysis says what in the repository most likely caused it. Keeping these two functions separate
makes the system easier to analyze and avoids overloading a single prompt with too many tasks at once.

\section{Stage 6: Patch Planning and Intended Action}

After root-cause analysis, PatchOps computes an intended action through a deterministic planner. This stage
combines failure category, confidence, and flaky status with the risk policy defined in Chapter 3. If the case
is high risk or unknown, the planner escalates. If it is medium risk, the planner defaults to preview. If it is
low risk and highly confident, the planner marks the case as a patch candidate. Otherwise, it remains
preview-oriented.

This stage is architecturally central because it is where repair becomes explicitly decision-theoretic. PatchOps
does not treat patch generation as the default outcome. Instead, it introduces a planner whose job is to decide
whether patch generation is even appropriate to pursue strongly.

\section{Stage 6b: Learned Decision Gate}

The learned decision gate is the implemented answer to the question of where the act-versus-hold threshold
comes from. Earlier versions of PatchOps used hand-set constants to decide whether a low-risk case was confident
enough for automatic patching. Phase 15 replaces that low-risk threshold test with an optional calibrated model
trained on the labeled Safety Set.

The gate uses only pre-generation signals: calibrated classifier confidence, category risk, flaky status,
semantic ambiguity, the number of files named in the log, and evidence length. Its output is a calibrated
probability of action. If the gate predicts act, the case may proceed toward patch generation and still must
pass the safety and plausibility gates. If it predicts hold, the case remains preview-oriented. The rest of the
policy is unchanged: flaky cases still rerun, high-risk and unknown cases still escalate, and low-plausibility
or unsafe patches are still downgraded later.

Architecturally, this means Stage 6b implements the Chapter 3 rule
\(
d_i = \textsc{act}
\)
iff
\(
p_{i,\mathrm{cal}} \ge 0.5
\)
on the low-risk lane only. It does not bypass the deterministic routing for flaky, medium-risk, high-risk, or
unknown cases, and it does not overrule the later safety and plausibility gates. This separation matters
because the learned model refines one bounded decision boundary inside the workflow rather than turning the
entire system into an opaque end-to-end controller.

This stage matters architecturally because it turns abstention from a threshold convention into an evaluated
model component. It also keeps the design conservative: the learned gate is opt-in and graceful. If the model
artifact is absent, the earlier deterministic threshold behavior remains available.

\begin{figure}[!htbp]
\centering
\includegraphics[width=\textwidth]{F20_decision_gate_architecture.jpg}
\caption[Learned decision gate architecture (Stage 6b)]{\textbf{Learned decision gate architecture (Stage 6b).}
Deterministic, zero-cost pre-generation signals
(calibrated classifier confidence, category-risk one-hot, flaky status, semantic ambiguity,
log-named-file count, and evidence length) are standardized and passed to a calibrated
logistic-regression model. The isotonic-calibrated probability of action $p_{\mathrm{cal}}$ is
thresholded at $0.5$ on the low-risk lane to set act (proceed to patch, then the safety and
plausibility gates) versus hold (preview for human review).}
\label{fig:decision-gate-arch}
\end{figure}

\section{Stage 7: Semantic Ambiguity Backstop}

Before any candidate automatic patch can survive the pipeline, the system applies a semantic ambiguity check.
This check is deterministic and looks for signals associated with type errors, assertion mismatches, and logic
exceptions. If such signals are present, an intended automatic patch is downgraded to preview.

Architecturally, this backstop is important because it enforces a boundary between mechanical and semantic
repair. It is also a good example of how PatchOps combines deterministic policy with LLM-driven reasoning:
models may propose a plausible repair path, but a separate policy layer can still decide that the action should
remain human-reviewed.

\section{Stage 8: Patch Generation}

When the planner and semantic guard still permit a code change, PatchOps generates a patch. This stage has both
a deterministic rule path and an LLM path. The deterministic path handles flagship mechanical cases such as
runtime version mismatches. The LLM path covers the broader set of repairable low- and medium-risk categories by
producing candidate file edits and then deriving a unified diff from them.

This hybrid architecture is significant for the thesis. It shows that PatchOps is not a pure LLM wrapper. Where
reliable deterministic repair is possible, the system prefers it. Where broader reasoning is required, it uses
an LLM under tightly constrained output expectations. The generated patch is accompanied by PR-facing metadata,
including title and explanation, which supports the system's live-review workflow.

\section{Stage 9: Safety Gate}

The safety gate is the core defensive component in the implemented architecture. It examines candidate patches
before they are allowed to remain in the automatic patch branch. The gate checks for forbidden paths,
category-specific allow-list violations, secret introduction, excessive patch size, and reward-hacking patterns
such as test tampering or check-disabling edits.

If the gate blocks a patch, the system never proceeds as though the change were safe for automatic use.
Depending on the case, the decision is downgraded to preview or escalate. This downgrade logic is architecturally
important because it shows that safety is not a separate report generated after the fact. It actively changes
the control flow of the system.

\section{Stage 10: Plausibility Gate}

Even a patch that passes the file-level safety gate may still be off-target. PatchOps therefore includes a
plausibility gate that estimates whether the generated patch is likely to correspond to a real fix or whether it
looks overfit. The plausibility score uses static signals such as evidence locality, minimality, category fit,
source-touch behavior, deletion heaviness, and reward-hack absence.

This stage is the architectural response to a central finding of the APR literature: plausible does not imply
correct. In PatchOps, a low-plausibility patch is not discarded silently. Instead, it is downgraded from
automatic patch to preview. This preserves usefulness while protecting the system from overclaiming confidence
in weak candidates.

\section{Final Decision and Trace}

After the previous stages have run, PatchOps emits one final decision. This decision, together with the
underlying evidence, reasoning trace, patch, and safety report, is returned as a structured result. The
trace-based design is intentional. Because the thesis emphasizes trustworthy automation, it is not enough for
the system to merely output an action. The action should be inspectable, explainable, and attributable to
intermediate stages.

This trace mechanism also supports the live demo interface and later analysis. A failed or downgraded patch is
not a black box; the user can inspect whether the downgrade happened because of low confidence, a safety-policy
violation, a semantic ambiguity signal, or low plausibility.

\section{Current Workflow as a Figure}

For thesis presentation, the current implemented system should be shown as a single-pass architecture diagram.
That figure starts from a CI failure log, passes through log extraction, classification, flaky detection,
repository retrieval, root-cause analysis, planning, patch generation, safety checking, and decision emission,
and ends in one of the four explicit actions. This is the correct ``implemented pipeline'' figure for the
thesis because it corresponds exactly to the built system rather than to a purely aspirational design.

\begin{figure}[!htbp]
\centering
\includegraphics[width=\textwidth]{F21_pipeline_overview.jpg}
\caption[The implemented PatchOps pipeline end-to-end]{\textbf{The implemented PatchOps pipeline end-to-end.}
A CI failure log is extracted and classified
(TF-IDF + calibrated logistic regression, then LLM adjudication); flaky failures rerun, otherwise the
system retrieves repository context, runs root-cause analysis, plans a risk class and intended action,
and sets act/hold at the decision gate (learned calibrated gate on the low-risk lane, hand-tuned
thresholds elsewhere). A semantic-ambiguity backstop, the safety gate, and the plausibility/overfit
gate then route each case into exactly one of four explicit decision outcomes: rerun, patch (open PR),
preview (human review), or escalate.}
\label{fig:pipeline-overview}
\end{figure}

\section{Validation Harness and Target Agentic Architecture}

Although the core decision engine is single-pass, PatchOps already includes a bounded local validation harness
used to apply candidate patches and re-run supported CI-style commands. This is the mechanism that makes
CI-validated repair measurable in the thesis. The stronger target architecture tightens this mechanism into two
additional capabilities: a verifier stage and a multi-attempt validation loop. These additions are what make the
architecture genuinely agentic in the stronger sense described earlier.

The verifier stage is an additional reasoning layer that critiques a generated patch before it reaches the final
safety and validation stages. Its purpose is to catch plausible-but-wrong fixes earlier. The validation loop
extends the workflow from one-shot repair into generate-validate-rediagnose behavior: if a patch is applied and
the CI-style validation still fails, the system re-enters diagnosis with the new error and attempts another
repair within a bounded retry limit.

Architecturally, this target design matters because it gives the thesis a clear progression path. The current
system is a carefully engineered decision pipeline with validation support. The target system is the same
pipeline extended with the smallest additional agentic pieces that the literature suggests are worth the
complexity. This keeps the design honest: agentic behavior is added where it promises measurable benefit, not
for rhetorical effect.

\section{Implemented Extensions and Target Growth}

The broader architecture also accommodates two important extensions that were introduced after the earliest base
pipeline and are now part of the implemented and evaluated PatchOps system: full-repository retrieval and a
lightweight learned classifier. Full-repository retrieval expands the candidate search space beyond log-named
files, addressing the localization bottleneck directly. The learned classifier supports reproducibility and
low-cost classification while preserving the main thesis focus on orchestration and safety.

It is important to separate these implemented extensions from the still-targeted verifier and validation-loop
additions described above. Full-repository retrieval and the learned classifier are no longer merely proposed;
they have already been built and measured in later project phases. Their role in this chapter is architectural:
they show that the pipeline is modular enough to accept stronger retrieval and classification components without
changing the underlying action model or safety policy. This modularity is one of the strongest design
properties of PatchOps as a research artifact.

\section{Why the Architecture Matches the Thesis Claim}

The PatchOps architecture is not merely a sequence of implementation modules. It is a direct structural
expression of the thesis claim. The system begins by compressing evidence rather than trusting raw logs,
retrieves bounded repository context rather than relying on generic code generation, separates diagnosis from
action choice, and places explicit policy barriers between a candidate patch and a final automatic action. In
doing so, it embodies four commitments:

\begin{enumerate}
\item \textbf{Repository awareness:} reasoning should be grounded in actual project files rather than logs alone.
\item \textbf{Safety gating:} candidate edits must satisfy deterministic policy checks before they are trusted.
\item \textbf{Abstention awareness:} the system should be able to choose non-patch actions as first-class
outcomes.
\item \textbf{Closed-loop validation:} stronger versions of the system should observe the consequences of a patch
through bounded re-execution rather than treating generation as the endpoint.
\end{enumerate}

These commitments are what distinguish PatchOps from both log-only CI assistants and unconstrained coding-agent
patterns.

\section{Chapter Conclusion}

This chapter described the PatchOps system design as a structured repair pipeline rather than a generic LLM
application. The architecture begins with deterministic evidence compression, moves through LLM-assisted
classification and root-cause analysis, and then narrows the decision through planning, the learned decision
gate, semantic ambiguity checks, safety gating, plausibility assessment, and final action selection. The current
engine is therefore a single-pass but decision-rich workflow, supported by a validation harness for selected
candidate fixes.

The chapter also distinguished this implemented architecture from the target agentic design, which adds a
verifier and a tighter bounded rediagnosis loop. That distinction is important for the thesis because it keeps
the implemented system and the proposed research direction clearly separated. With the architecture now
established, the next chapter can map these stages to the actual codebase and implementation details.
